\section*{Rare shared-coordinate degree-shift counterfamily}

Put
\[
 \phi(t)=2t^2\left(\log_2\frac1t+\frac1{\ln2}\right),\qquad
 \Phi(F)=\sum_j\phi(|\widehat F(\{j\})|).
\]
The accepted disjoint-variable argument uses
$\Phi(BF)\leq \Phi(F)+M I(B)$.  The following genuine Boolean family
shows that no shared-variable version with a bounded charge on the shared
physical edge can hold.

Let $n=2m+1$, let $K_n$ be majority on $n$ signs $u_1,\ldots,u_n$, and
write
\[
 c_m=\widehat K_n(\{u_i\})={\binom{2m}{m}\over 4^m}.
\]
Let $x$ be one further sign, let $v=(v_1,\ldots,v_r)$, let $E(v)$ be one
fixed point of the $r$-cube, and put $\epsilon=2^{-r}$.  Define
\[
 f(x,u)=xK_n(u),\qquad
 g(x,v)=\begin{cases}x,&E(v),\\1,&\neg E(v),\end{cases}\qquad h=fg.
\]
All three functions are sign-valued.  The two factors share only the
physical coordinate $x$.

\subsection*{Exact Fourier and influence ledger}
Balancedness of $K_n$ gives
\[
 \Phi(f)=0,qquad \Phi(g)=(r+1)\phi(\epsilon),qquad
 \Phi(h)=n\phi(\epsilon c_m).
\]
Indeed multiplication by $x$ moves every singleton coefficient of $K_n$
to the pair $\{x,u_i\}$, while the rare factor moves an $\epsilon$ fraction
back to $\{u_i\}$.  The relevant physical influences are
\[
 \operatorname{Inf}_x(f)=1,quad \operatorname{Inf}_x(g)=\epsilon,quad
 \operatorname{Inf}_x(h)=1-\epsilon,quad I(g)=(r+1)\epsilon.
\]
Thus the exact overlap saving on the shared edge is $2\epsilon$.

For completeness, if two Boolean factors share only $x$ and have independent
private variables, writing $D_xF=(F_+-F_-)/2$ and
$\bar F=(F_++F_-)/2$, the pointwise Boolean identity
$D_xF\,\bar F=0$ gives
\[
 \operatorname{Inf}_x(FG)=p(1-q)+(1-p)q=p+q-2pq.
\]
For a private coordinate $u_i$ of $F$ the singleton convolution formula is
\[
 \widehat{FG}(\{u_i\})
 =\widehat F(\{u_i\})\widehat G(\emptyset)
  +\widehat F(\{u_i,x\})\widehat G(\{x\}).
\]
The second term is the degree-shift mechanism absent in the disjoint proof.

\subsection*{Rigorous divergence}
The elementary induction bounds
\[
 {1\over2\sqrt m}\le c_m\le {1\over\sqrt{m+1}}
\]
follow from $c_{m+1}/c_m=(2m+1)/(2m+2)$.  Consequently
\[
 \Phi(h)\ge {\epsilon^2\over2}\log_2(m+1).
\]
Therefore, for every $r\ge1$ and every finite $C$, every $m$ satisfying
\[
 {\epsilon^2\over2}\log_2(m+1)>
 (r+1)\phi(\epsilon)+C(1+\epsilon)
\]
violates
\[
 \Phi(fg)\le\Phi(f)+\Phi(g)+C
 \bigl(\operatorname{Inf}_x(f)+\operatorname{Inf}_x(g)\bigr).
\]
The same family violates $\Phi(fg)\le\Phi(f)+M I(g)$ for every fixed
finite $M$ by replacing the last term by $M(r+1)\epsilon$.  In particular,
neither the rare first-revelation pivotal mass nor the total shared-edge mass
can pay the newly exposed singleton potential with an absolute multiplier.

\subsection*{Reveal-order comparison}
This is a local product-law obstruction, not a restriction-tree no-go.
After fully revealing the controls $v$ but not $x$, the expected singleton
potential is exactly
\[
 \epsilon\Phi(K_n).
\]
If instead $x$ is revealed first, the two children are $K_n$ and
$K_nQ_r$, where $Q_r=2\mathbf1_E-1$ and
$|\mathbb E Q_r|=1-2\epsilon$.  Hence the expected intermediate potential is
\[
 {n\over2}\{\phi(c_m)+\phi((1-2\epsilon)c_m)\}.
\]
After both $x$ and all controls are revealed, the terminal potential is
always $\Phi(K_n)$.  Thus the total signed potential telescope is
order-independent, but the location of its large negative refund is not.
Any viable physical stopping ledger must carry this signed degree-shift
balance across nodes; a nonnegative shared-pivotal account loses it.

This family does not refute the exact mismatch/KL identity, existential
coordinate selection, an optimized globally chosen tree, or FEI.  It only
refutes the stated overlap-stable singleton-potential/product payments.
